\chapter{Das Und}

Wer \zitat{Zeit und Raum} sagt, hat schon entschieden, ohne zu wissen wie. Die Konjunktion stellt zwei Namen nebeneinander, als wären sie Nachbarn. Aber Nachbarschaft ist eine räumliche Beziehung. Sagt man stattdessen, Zeit und Raum seien \emph{zugleich} da, so hat man die Zeit zur Bühne gemacht, auf der beide stehen. Eugen Fink hat diese Verlegenheit als erste Frage einer Untersuchung von Raum, Zeit und Bewegung festgehalten: Wie sind die beiden beisammen? Nebeneinander oder gleichzeitig? Keine der beiden Antworten ist neutral, jede nimmt Partei für eine der beiden Seiten, und zwar bevor irgendetwas über die Sache gesagt ist.\footnote{Fink, Zur ontologischen Frühgeschichte von Raum -- Zeit -- Bewegung, Vorverständigung.}

Julius Baumann hat die Parteinahme, die in der Geschichte des Denkens tatsächlich stattgefunden hat, in einem Satz zusammengefasst: \zitat{Die Zeit hat das Schicksal gehabt, gewöhnlich nach dem Raum und als diesem parallel behandelt zu werden.}\footnote{Baumann, Die Lehren von Raum, Zeit und Mathematik, Bd.\,2, Schluss.} Das Wort \zitat{Schicksal} ist genau. Es war keine Entscheidung, die jemand getroffen hätte, sondern eine Gewohnheit, die sich aus der Sache selbst zu ergeben schien: Der Raum liegt vor, er lässt sich zeichnen, messen, teilen und wieder zusammensetzen. Die Zeit liegt nicht vor. Was von ihr da ist, ist ein Jetzt, das schon nicht mehr ist, wenn man es zeigt. Also nimmt man das Vorliegende als Muster für das Nichtvorliegende: eine Linie für die Zeit, Punkte für die Augenblicke, Strecken für die Dauern. Und weil das Muster so gut passt, bemerkt man nicht, was es verdeckt.

Es verdeckt drei Dinge. Es verdeckt, dass die Zeit ihre Grenze anders hat als der Raum: Der Raum ist als Ganzes da und lässt sich überall begrenzen, die Zeit ist nur als Grenze da und lässt sich nirgends beliebig begrenzen. Es verdeckt, dass die Zeit eine Richtung hat, die kein Raum aus sich hervorbringt; wo der Raum Richtungen zeigt, oben und unten, vorn und hinten, rechts und links, stammen sie von einem Leib, der in ihm lebt, und der Leib ist ein Zeitliches. Und es verdeckt, dass das Maß, mit dem der Raum gemessen wird, aus der Zeit kommt: Das Lichtjahr ist keine astronomische Bequemlichkeit, sondern die Form aller Längenmessung, sobald man fragt, was Länge physikalisch bedeutet.

Grenze, Richtung, Maß: Das sind die drei Asymmetrien, um die es im Folgenden geht. Sie werden nicht behauptet, sondern an den Texten aufgesucht, die von der Sache handeln. Diese Texte gehören zwei Lagern an. Das eine Lager ist das der Parallelisierung. Ihm gehören Zenon an, wie Fink ihn liest, Kant, der beide Anschauungsformen in einem Kapitel behandelt, Minkowski, der die Zeit zur weiteren Koordinate macht, und Hegel, der Raum und Zeit \zitat{zwei Weisen des Außersichseins} nennt. Das andere Lager ist das der Asymmetrie: Reichenbach, für den die Zeit die Faserrichtung der Kausalketten ist und der Raum nur ihr Nebeneinander; Brentano, für den das räumliche Kontinuum als Ganzes besteht und das zeitliche nur als Grenze; Guzzoni, für die die Zeit zum Raum wird und nicht umgekehrt; Koch, für den kein Raumpunkt von seinem Spiegelbild unterscheidbar ist, solange nicht ein Leib das Hier besetzt. Und Hegel gehört auch dem zweiten Lager an, denn bei ihm geht der Raum in die Zeit über, und \zitat{die Zeit ist also erst der wahrhafte Punkt}.\footnote{Hegel, Vorlesung über Naturphilosophie, Nachschrift Uexküll, zitiert nach Bonsiepen, Hegels Raum-Zeit-Lehre.}

Dass Hegel in beiden Lagern steht, ist kein Zufall und keine Inkonsequenz. Es zeigt, was die Parallelisierung eigentlich ist: nicht ein Irrtum, sondern eine Abstraktion, die ihre eigene Aufhebung schon enthält. Wer Zeit und Raum parallel behandelt, muss an irgendeiner Stelle sagen, wodurch sie sich unterscheiden. Er legt diese Stelle nur dorthin, wo sie am wenigsten auffällt: in ein Vorzeichen, in eine Fußnote, in einen Übergang. Die Aufgabe ist deshalb nicht, die Parallelisierung zu widerlegen, sondern die Stellen aufzusuchen, an denen sie die Asymmetrie versteckt, und sie dort herauszuholen.

Was dabei herauskommt, ist keine Umkehrung des Schicksals. Es wäre leicht, nun den Raum nach der Zeit und als dieser parallel zu behandeln, und einige der besten Gedanken, die im Folgenden vorkommen, sind in Gefahr, genau das zu tun: Reichenbachs Reduktion der Raummessung auf Zeitmessung, Hegels Übergang des Raumes in die Zeit, Heideggers frühe Rückführung der Räumlichkeit auf die Zeitlichkeit. Aber jeder dieser Gedanken trägt seinen eigenen Vorbehalt bei sich. Reichenbach braucht einen materiellen Messkörper, den die Lichtgeometrie nicht liefert. Kant braucht ein Beharrliches im Raum, ohne das die Zeit nicht einmal als Größe vorgestellt werden könnte. Heidegger hat seine Rückführung zurückgenommen: Sie \zitat{läßt sich nicht halten}.\footnote{Heidegger, Zeit und Sein, zitiert nach Guzzoni, Weile und Weite.} Der Raum widersteht. Was er ist, wenn er nicht das Parallele der Zeit und nicht ihr Abgeleitetes ist, bleibt am Ende offen, und es wird nicht so getan, als wäre es geschlossen.
